\section{Introduction}\label{sec:introduction}

The absolute values of unramified Hecke eigenvalues measure the departure
of an automorphic representation from temperedness. The Ramanujan--Arthur
prediction organizes the departures that occur in the discrete spectrum
by homomorphisms from \(\mathrm{SL}_2\) into the Langlands dual group;
the additional algebraic factor is central to Arthur's conjectural
description of non-tempered discrete automorphic representations
\cite[Sections~4 and~8]{Arthur1989}.
For cuspidal automorphic functions over a global function field, this leads to a decomposition
indexed by nilpotent orbits. We prove that decomposition with arbitrary
full finite level. Its common orbit also turns genericity at one
unramified place into temperedness at all unramified places in the
adjoint absolutely simple case.

\subsection{The statement}

Let \(X\) be a smooth projective geometrically connected curve over
\(\mathbb F_q\), let \(F=\mathbb F_q(X)\), and let \(G\) be a split
connected semisimple group over \(\mathbb F_q\). Write \(\mathbb A_F\)
for the adele ring and \(\mathcal O_{\mathbb A}\) for its integral
subring. Let \(D\) be any effective divisor on \(X\) defined over
\(\mathbb F_q\); its multiplicities are arbitrary. Set
\begin{equation}\label{intro:level}
 K_D=\ker\bigl(G(\mathcal O_{\mathbb A})\longrightarrow
                   G(\mathcal O_D)\bigr),
 \qquad
 \mathcal X_D=G(F)\backslash G(\mathbb A_F)/K_D.
\end{equation}
For \(D=0\), this means \(K_D=G(\mathcal O_{\mathbb A})\).
Equivalently, \(\mathcal X_D\) is the set of isomorphism classes in
\(\Bun_{G,D}(\mathbb F_q)\), where the level is a trivialization on
the whole finite subscheme \(D\).

Let \(C_{D,\mathbb Q}\) be the space of compactly supported
\(\mathbb Q\)-valued functions on \(\mathcal X_D\) satisfying
\begin{equation}\label{intro:cuspidality}
 \int_{U(F)\backslash U(\mathbb A_F)} f(ug)\,du=0
\end{equation}
for every proper \(F\)-parabolic subgroup \(P\subset G\), its
unipotent radical \(U\), and every \(g\in G(\mathbb A_F)\).
This condition is independent of the nonzero normalization of \(du\).
For a prime \(\ell\ne\operatorname{char}(\mathbb F_q)\), put
\[
 C_{D,\ell}=\overline{\mathbb Q}_\ell\otimes_{\mathbb Q}C_{D,\mathbb Q}.
\]
All these cusp spaces at fixed level are finite dimensional.

Fix the pinned split rational form of the dual group \(\Gd\), with Lie
algebra \(\gd\), and the rational conventions of
\cite[Sections~3.1.1 and 3.4.1--3.4.3]{GLR}.
For a nilpotent \(\Gd\)-orbit \(\mathcal O\subset\gd\), choose a
Jacobson--Morozov homomorphism
\[
 \phi_{\mathcal O}:\mathrm{SL}_2\longrightarrow\Gd,
 \qquad
 H_{\mathcal O}=Z_{\Gd}\bigl(\phi_{\mathcal O}(\mathrm{SL}_2)\bigr).
\]
The centralizer \(H_{\mathcal O}\) may be disconnected. Define
\(H_{\mathcal O}^{c}(\overline{\mathbb Q})\) to consist of its
semisimple algebraic points whose images under every embedding
\(\overline{\mathbb Q}\hookrightarrow\mathbb C\) are conjugate
into a maximal compact subgroup of \(H_{\mathcal O}(\mathbb C)\).
For a prime power \(r\), write
\begin{equation}\label{intro:orbit-spectrum}
 \begin{split}
 r_{\mathcal O}
  &=\phi_{\mathcal O}
       \begin{pmatrix}r^{1/2}&0\\0&r^{-1/2}\end{pmatrix},\\
 R_{r,\mathcal O}
  &=\operatorname{image}\left(
       r_{\mathcal O}H_{\mathcal O}^{c}(\overline{\mathbb Q})
       \longrightarrow
       (\Gd/\!/\operatorname{Ad}(\Gd))(\overline{\mathbb Q})
     \right).
 \end{split}
\end{equation}
We fix compatible square roots and use the rational Satake normalization
of \cite[Equation~(3.2)]{GLR}. Changing a square root multiplies the displayed
factor by \(\phi_{\mathcal O}(-I)\), a finite-order element of
\(H_{\mathcal O}\), and does not change the set
in Equation~\eqref{intro:orbit-spectrum}.

For a closed point \(x\notin\operatorname{supp}(D)\), put
\(d_x=[k(x):\mathbb F_q]\) and \(q_x=q^{d_x}\). Its spherical
Hecke algebra is
\[
 A_x=\operatorname{Funct}_c
       \bigl(G(\mathcal O_x)\backslash G(F_x)/G(\mathcal O_x),\mathbb Q\bigr),
\]
with Haar measure normalized by
\(\operatorname{vol}(G(\mathcal O_x))=1\). We regard each character
of \(A_x\) over \(\overline{\mathbb Q}\) as a semisimple dual-group
conjugacy class by this fixed Satake convention.

Define \(C_{D,\mathcal O,\mathbb Q}\subset C_{D,\mathbb Q}\) by the
following spectral-support condition: for every
\(x\notin\operatorname{supp}(D)\), every character occurring in
\[
 \overline{\mathbb Q}\otimes_{\mathbb Q} A_xf
\]
has Satake class in \(R_{q_x,\mathcal O}\). Set
\(C_{D,\mathcal O,\ell}
 =\overline{\mathbb Q}_\ell\otimes_{\mathbb Q}
 C_{D,\mathcal O,\mathbb Q}\).
The same orbit is required at all good points.

\begin{theorem}\label{main:decomposition}
For every \(X,G,D,\ell\) as above, the natural addition maps
\begin{equation}\label{intro:decomposition}
 \bigoplus_{\mathcal O}C_{D,\mathcal O,\mathbb Q}
       \xrightarrow{\ \sim\ }C_{D,\mathbb Q},
 \qquad
 \bigoplus_{\mathcal O}C_{D,\mathcal O,\ell}
       \xrightarrow{\ \sim\ }C_{D,\ell}
\end{equation}
are isomorphisms. The sums range over the nilpotent orbits of \(\Gd\).
\end{theorem}

Taking \(D=0\) resolves Conjecture~3.4.6 of
Gaitsgory--Lafforgue--Raskin positively, as well as its equivalent rational
form, Conjecture~3.4.5. Theorem~\ref{main:decomposition} also establishes
the stated finite-level extension. The compactness in
Equation~\eqref{intro:orbit-spectrum} is the archimedean condition at
every embedding; no extra nonarchimedean condition is included.

\subsection{Generic cuspidal representations}

For a split group, fix a Borel subgroup with unipotent radical \(N\).
A local representation of \(G(F_x)\) is \emph{generic} if it admits a
nonzero Whittaker functional for a nondegenerate character of \(N(F_x)\).
A cuspidal automorphic representation is \emph{globally generic} if one
of its vectors has a nonzero global Whittaker coefficient for a
nondegenerate character of \(N(F)\backslash N(\mathbb A_F)\).
The generalized Ramanujan conjecture predicts that globally generic
cuspidal representations are tempered at every place
\cite[Section~1.1]{CiubotaruHarris}. We obtain its unramified assertion
in the following scope, with genericity required at only one unramified
place.

\begin{corollary}[Unramified generalized Ramanujan]
\label{main:generic-ramanujan}
Let \(F\) be the function field of a smooth projective geometrically
connected curve over a finite field, and let \(G/F\) be split connected
adjoint absolutely simple. Let
\(\pi=\bigotimes'_x\pi_x\) be a complex cuspidal automorphic
representation of \(G(\mathbb A_F)\). If \(\pi_v\) is unramified and
generic at one place \(v\), then \(\pi_x\) is tempered at every place
\(x\) at which it is unramified. In particular, every globally generic
cuspidal automorphic representation of \(G(\mathbb A_F)\) is tempered
at every unramified place.
\end{corollary}

Here unramified means spherical for a hyperspecial maximal compact
subgroup. The conclusion is only about these local components. The
proof, given at the end of Section~\ref{sec:global}, combines the common
orbit with the local geometric-parameter criterion of Ciubotaru and
Harris.

\subsection{Context}

The Langlands correspondence for general linear groups over function
fields, established by Drinfeld in rank two and by Laurent Lafforgue
in general, gives the purity theorem used at the end of this proof
\cite[Introduction and Theorem~VII.6]{LLafforgue}.
Vincent Lafforgue's construction of
global parameters extends the parameterization of cuspidal functions to
split reductive groups with arbitrary finite level
\cite[Theorem~1.1]{VLafforgue}. These results provide global parameters
and control the weights of their general-linear constituents. The main
issue here is to constrain the absolute-value part of an unramified
dual-group parameter to the neutral cocharacter of a nilpotent orbit.

Gaitsgory--Lafforgue--Raskin connect nilpotent-orbit filtrations on
automorphic functions with Hecke spectral support and formulate the
rational and \(\ell\)-adic cusp decompositions in
\cite[Section~3.4]{GLR}. Their Remark~3.4.7 identifies the latter as
the \(\ell\)-adic Ramanujan--Arthur conjecture. We address this
spectral-support assertion directly.

Using the present decomposition at \(D=0\), the companion manuscript
identifies each cuspidal canonical
Arthur-filtration step with the scalar extension of the corresponding sum
of the rational summands here, under its four characteristic hypotheses
\cite[Assumption~1.1 and Corollary~7.5]{ArthurRationality}. This is the
common-scope form of \cite[Conjectures~3.4.9 and~3.4.11]{GLR}.

Earlier function-field results constrain the unramified spectrum under
additional local hypotheses. In characteristic greater than two, Sawin and
Templier prove unramified temperedness under monomial geometric
supercuspidality and base-change assumptions
\cite[Theorem~1.1]{SawinTemplier}; their Theorem~11.7 also shows that a
cuspidal representation of a split semisimple group tempered at one
unramified place is tempered at all such places. Ciubotaru and Harris
obtain restrictions to complementary series indexed by nilpotent orbits,
and generalized Ramanujan under additional local assumptions, from
excursion parameters, Frobenius weights, and the spherical unitary dual
\cite[Theorems~1.11 and~5.4]{CiubotaruHarris}.
Theorem~\ref{main:decomposition} supplies the nilpotent-orbit restriction
for the whole cuspidal space at full finite level.
Corollary~\ref{main:generic-ramanujan} uses their local criterion
\cite[Corollary~6.32\textup{(1)}]{CiubotaruHarris}: the algebraic
\(\mathrm{SL}_2\) cocharacter in our decomposition makes the Satake
parameter geometric, and the
common orbit propagates the resulting temperedness without the additional
tempered-place hypothesis of their global Theorem~5.4.

The geometric ingredients have a separate ancestry. Geometric Satake
identifies spherical sheaves with representations of the dual group
\cite{MV}; its derived form and the regular-sheaf and Springer
calculations are due to Ginzburg, Bezrukavnikov--Finkelberg, and
Arkhipov--Bezrukavnikov--Ginzburg \cite{GinzburgLoop,BF,ABG}.
We use the enhanced arithmetic Satake formulation of
Ben-Zvi--Sakellaridis--Venkatesh \cite{BZSV} together with the
tensor-product and trace formalism of Ben-Zvi--Francis--Nadler
\cite{BFN}.
Gaitsgory's nearby-cycle construction supplies central sheaves
\cite{GaitsgoryCentral}, and the Wakimoto filtration and its central
compatibilities are developed by Arkhipov--Bezrukavnikov and, in the
integral parahoric setting, Cass--van den Hove--Scholbach
\cite{AB,CvdHS}. Section~\ref{sec:categorical} states the precise
versions used here and proves the additional specialization,
root-datum, and cyclic-transfer compatibilities needed for the argument.

\subsection{Proof strategy}

The main intermediate result is local. An ordered Bernstein weight is
a joint character of the commuting Iwahori translation operators,
before passage to the Weyl-group quotient. If its absolute-value
exponent is dominant and it occurs in the discrete automorphic
spectrum, then it is a compact
factor times a Jacobson--Morozov value
(Theorem~\ref{ext:discrete}). We prove this by contradiction and
induction on semisimple rank. For a hypothetical violating weight \(t\)
at \(x\), choose a complex realization and write
\(\lambda=\log_{q_x}|t|\) in the dominant chamber; unitarity makes
its central component zero. The contradiction comes from two degree bounds on
the same spherical cohomological summand.

First, standard translations define bounded open shtukas that detect
the prescribed ordered weight. A polynomial filter removes irrelevant
continuous spectral families; induction gives a strict spectral gap
for those that remain. The rotation trace identity then forces the
weight to occur in compact cohomology. Deligne's weight bound puts its
degree at least
\(2d_x\langle\lambda,\mu\rangle-C'\), where \(\mu\) is the
translation weight and \(C'\) is independent of \(\mu\).
The occurrence itself, through integrality of Frobenius weights, also
forces \(\lambda\) to be rational
(Corollary~\ref{ext:rationality}). A fixed rational Weyl-invariant
metric therefore lets us choose integral highest weights
\(\mu=m\lambda\) for sufficiently divisible positive integers \(m\),
including when \(\lambda\) lies on a chamber wall.

The Wakimoto filtration compares this translation cohomology with
cohomology for the spherical Satake label of highest weight \(\mu\).
Since \(\mu=m\lambda\) is the unique weight maximizing pairing
with \(\lambda\), its extreme part cannot cancel against another
filtration term. The hyperspecial projections retain that part at a
spherical level. There the categorical trace realizes the projection
on complexes, and a finite specialization argument preserves its top
nonzero cohomology. This produces the spherical summand on which both
degree bounds can be compared.

The upper bound is prepared independently in
Section~\ref{sec:amplitude}. Derived Satake and the categorical trace
calculation express the spherical trace as an equivariant module.
Its group-valued trace variable \(s\in\Gd\) and Lie-algebra variable
\(e\in\gd\) satisfy the derived relation
\begin{equation}\label{intro:resonance}
 \operatorname{Ad}(s)e=q_xe.
\end{equation}
At semisimple part \(t\), the resonant space is
\(E_t=\{e\in\gd:\operatorname{Ad}(t)e=q_xe\}\).
Uniform bounds from costandard translations and nef Springer line
bundles pass, by Levi restriction and localization on opposite flags,
to bounds on whole derived fibers, with degrees corrected by compatible
\(\mathfrak{sl}_2\)-triples. A finite-orbit induction then
gives the upper bound for equivariant modules satisfying these tests,
without a coherence hypothesis. Applied to the summand just constructed, the two
bounds on its top nonzero degree \(N\) are
\[
 2d_x\langle\lambda,\mu\rangle-C'
 \ \leq\ N\ \leq\
 C+\max_{e\in E_t}d_x\langle\mu,h_e\rangle,
\]
where \(h_e\) is the neutral cocharacter of a compatible triple and
\(C,C'\) are independent of \(m\). On \(\mu=m\lambda\), the
difference of the leading terms is
\[
 2d_xm\min_{e\in E_t}\|\lambda-h_e/2\|^2.
\]
It is positive for a violating parameter, contradicting the two
bounds as \(m\) grows. This proves the local theorem.

Two parts of the argument can be used separately. Propositions
\ref{amp:levi} and \ref{amp:fiber} extract corrected fiber bounds for
arbitrary equivariant modules from uniform flag tests.
Lemma~\ref{rot:picard} compares two descent conditions
by realizing a divisible rotation in a connected generalized Picard
group and applying Lang's theorem. The latter construction keeps the
entire level modulus and treats arithmetic components of a splitting
cover explicitly.

Finally, exact root equations for compatible triples show that the
compact factor can be chosen algebraically and is compact at every
complex embedding. Purity of the irreducible constituents of the global
parameter makes the dominant neutral cocharacter independent of the good place.
The finite rational Hecke algebra then gives the direct sum in
Equation~\eqref{intro:decomposition}.

\subsection{Organization and conventions}

Section~\ref{sec:categorical} constructs the cyclic trace and records
its Satake, central, and Springer descriptions.
Section~\ref{sec:amplitude} proves the amplitude estimate.
Section~\ref{sec:rotation} proves the rotation trace identity with
full finite level.
Section~\ref{sec:extraction} extracts the local constraint from the
discrete spectrum.
Section~\ref{sec:global} proves Theorem~\ref{main:decomposition} and
deduces Corollary~\ref{main:generic-ramanujan}.

The local arguments are formulated for split connected reductive groups
so that they apply inductively to Levi subgroups. Central degree is
bounded modulo a discrete lattice where necessary. On the sheaf side
the coefficients have characteristic zero, and all derived tensor
products and fibers are retained. A complex realization is chosen only
when absolute values are used. At a place of degree \(d\), the Hecke
parameter is \(q^d\), while the global cohomological Frobenius has base
\(q\). Standard translations are used for the trace identity;
costandard translations are used for upper amplitude estimates.
The established sheaf-theoretic and spectral inputs are stated with
their precise roles in the sections where they enter.
